\section{Proposed Methodology}
\label{sec:method}
\subsection{Overall Framework}
\label{sec:framework}
The proposed Diff-LM-GZSL framework aims to address the challenges of zero-shot fault diagnosis (ZSFD) in industrial systems by establishing a continuous semantic alignment between high-dimensional vibration data and expert-level natural language descriptions. Traditional ZSFD methods often rely on discrete attribute vectors or simple class names, which fail to capture the complex, nuanced characteristics of diverse fault modes. As illustrated in Fig. \ref{fig:framework}, our framework leverages the generative capabilities of latent diffusion models (LDMs) guided by language-driven semantic embeddings to synthesize robust visual features for unseen fault categories. This innovative approach recognizes that the mapping between mechanical health states and their corresponding vibration signatures is inherently stochastic and significantly dependent on the operational context. By integrating deep generative modeling with large-scale linguistic knowledge, Diff-LM-GZSL provides a versatile solution that can adapt to specialized machinery with minimal domain-specific fine-tuning.

The architecture comprises three primary modules: (1) Language-Driven Semantic Embedding, (2) Conditional Latent Diffusion Model (CLDM), and (3) Semantic Alignment and Classification. The Language-Driven Semantic Embedding module transforms raw fault descriptions—curated from expert knowledge—into a continuous, high-dimensional semantic space using a pre-trained large language model (LLM). This ensures that the global relationship between known and unknown faults is preserved through shared linguistic concepts. The CLDM then utilizes these embeddings as conditioning signals to model the conditional distribution of vibration features. By operating in a compressed latent space, the diffusion process effectively captures the underlying manifold of mechanical health states while mitigating the computational burden associated with raw signal generation. Finally, the Semantic Alignment and Classification module ensures that the generated features and real features are mapped back to a joint embedding space where categorical boundaries are strictly enforced. This three-stage modular approach allows for independent optimization of each component, ensuring that signal processing, semantic understanding, and generative synthesis are all fine-tuned for the task.

The operational workflow of Diff-LM-GZSL is divided into a training phase and an inference (testing) phase. During training, the model is optimized using data from seen fault classes. The feature extractor learns to map raw signals to a discriminative feature space, while the CLDM is trained to reconstruct these features from noise, conditioned on the corresponding language-derived semantic vectors. This training process involves a forward diffusion phase where Gaussian noise is incrementally added to the latent features, and a reverse denoising phase where the model learns to recover the original feature under the guidance of class-specific text embeddings. In the inference stage, for a fault class that was never encountered during training (an unseen class), we first construct its natural language description and project it into the semantic space. The trained CLDM then synthesizes a set of "pseudo-features" for this unseen class. These synthesized features, along with real features from seen classes, are used to train a final softmax classifier or a distance-based metric learner, enabling the model to perform generalized zero-shot fault diagnosis (GZSL) across the entire label space. This dual-phase operation effectively bridges the gap between historical training data and the unpredictable nature of real-world machine failures. As shown in IV-A, the overall architecture is visualized in Fig.~\ref{fig:framework}.

\subsection{Data Preprocessing and Feature Extraction}
\label{sec:preprocessing}
To handle the high-frequency and non-stationary nature of industrial vibration signals, a rigorous data preprocessing and feature extraction pipeline is implemented. Raw signals collected from accelerometers or pressure sensors are typically contaminated with environmental noise and operating condition variations. We employ a sliding window strategy to segment the continuous time-series data into fixed-length frames, ensuring sufficient temporal resolution for fault signature detection. This segmentation balances the trade-off between capturing transient anomalies and maintaining structural continuity within the data stream. Each segment is normalized using Z-score standardization to ensure zero mean and unit variance, which stabilizes the convergence of the downstream neural networks and prevents any single feature from dominating the learning process due to differences in scale. Moreover, signal processing techniques such as the Fast Fourier Transform (FFT) or Wavelet transforms can be integrated to refine the signal's spectral properties prior to feature learning.

The feature extraction network is designed as a deep 1D Convolutional Neural Network (1D-CNN) or a Stacked Autoencoder (SAE), depending on the complexity of the input data and the specific requirements of the industrial monitoring system. In this study, we focus on a 1D-CNN architecture due to its superior ability to capture local shift-invariant patterns in raw vibration waveforms without the need for manual feature engineering. These deep models (SAE/1D-CNN) have been widely proven effective in mechanical fault feature extraction \cite{Vincent2010SAE, LeCun2015CNN}. The network consists of multiple convolutional layers with varying kernel sizes to extract multi-scale spatio-temporal features, which are vital for characterizing complex phenomena like bearing spalls or gear pitting. These layers are meticulously followed by batch normalization and Rectified Linear Unit (ReLU) activation functions to introduce non-linearity and accelerate training. Global average pooling is applied at the final layer to reduce the spatial dimensions and aggregate the learned representations into a compact, discriminative signature.

Formally, let the $i$-th input signal segment be $x_{raw}^{(i)} \in \mathbb{R}^L$, where $L$ is the window length. The signal is mapped to a low-dimensional feature vector $x^{(i)} \in \mathbb{R}^d$ through the composite nonlinear mapping function $f_{\phi}(\cdot)$:
\begin{equation}
x = f_{\phi}(x_{raw})
\end{equation}
where $\phi$ represents the learnable parameters of the deep feature extractor, encompassing weights and biases of the convolutional and pooling layers. The formatting of the input data emphasizes the preservation of phase coherence and frequency distribution, which is critical for identifying specific fault types in high-speed rotating machinery. Detailed configurations of the feature extraction network, including the specific number of filters, kernel sizes, and stride lengths for each layer, are summarized in Table \ref{tab:network} and referenced throughout the evaluation. This compact representation $x$ serves as the primary target for the conditional latent diffusion process, effectively bridging the raw sensory input and the latent semantic space. By decoupling the feature extraction from the generative task, our framework can leverage pre-trained diagnostic models as foundations for zero-shot expansion.

\subsection{Language-Driven Semantic Embedding}
\label{sec:semantic}
The core innovation of the Diff-LM-GZSL framework lies in the substitution of rigid, manually-defined binary attribute tables with a dynamic, language-driven semantic embedding. In industrial fault diagnosis, fault characteristics are often best described using the qualitative language of field maintenance experts. We construct a comprehensive fault knowledge base where each fault class $y$ is associated with a detailed natural language description $T_y$. These descriptions encompass spectral characteristics, underlying physical phenomena, and typical temporal evolution patterns for each mechanical health state. For instance, an "Outer race fault" might be described as: "Outer race fault characterized by high-frequency periodic impulses in the vibration signal, accompanied by significantly increased root mean square (RMS) amplitude and harmonics of the ball pass frequency on the outer race (BPFO)." This text-based grounding provides a much more flexible and descriptive basis for classification than traditional one-hot encodings or fixed attribute lists.

To convert these rich textual descriptions into consistent mathematical representations, we utilize a pre-trained Transformer-based language encoder (e.g., BERT \cite{Devlin2019BERT}, CLIP \cite{Radford2021CLIP}, or a specialized LLM adapted for technical text). Let $T$ represent the string of the fault description. The encoding process extracts the dense semantic information from the text:
\begin{equation}
s = \text{Encoder}_{\text{LLM}}(T)
\end{equation}
where $s \in \mathbb{R}^{d_{LLM}}$ is the raw semantic vector extracted from the final embedding layer, typically by taking the [CLS] token or via mean-pooling across the sentence's sequence. Since the native dimensionality $d_{LLM}$ may not match the conditioning requirements of the latent diffusion model, we introduce a learnable semantic projection layer:
\begin{equation}
a = W_s s + b_s
\end{equation}
where $W_s \in \mathbb{R}^{d_s \times d_{LLM}}$ is a projection matrix and $b_s \in \mathbb{R}^{d_s}$ is the bias term. The resulting vector $a \in \mathbb{R}^{d_s}$ constitutes the final semantic embedding that encapsulates the diagnostic essence of the fault in a continuous manifold.

Compared to traditional binary attributed-based zero-shot methods, our language-driven approach offers several distinct advantages. First, it provides a significantly higher information density; a single sentence can describe the complex interaction between different mechanical components (e.g., "gear pitting with sequential lubrication failure"), whereas binary attributes require a pre-defined and often limited vocabulary. Second, the continuous nature of the LLM-derived semantic space allows the model to capture "near-miss" similarities between faults that might share linguistic descriptions but exist in different regions of the feature space. This inherent continuity is vital for generalizing to unseen fault categories, as it provides a smoother, more interpretable manifold for the diffusion process to interpolate across. Furthermore, this method eliminates the time-consuming and subjective process of manual attribute labeling. The contrast between traditional binary or discrete attributes and our proposed holistic text descriptions is further detailed in Table \ref{tab:network} and visualized in the detailed semantic module architecture in Fig.~\ref{fig:semantic_module}. By grounding the fault diagnosis in natural language, we effectively transfer broad human expert knowledge into the deep learning model, significantly enhancing its semantic reasoning capabilities and enabling true zero-shot performance in evolving industrial environments. Through Eq. 5 and Eq. 6, we establish a robust pipeline that can ingest arbitrary technical descriptions and convert them into potent diagnostic signals.